\chapter{L'intégration continue (CI)}

Un article en \LaTeX{} qui compilait hier ne compile plus aujourd'hui ; une preuve formelle
qui passait se casse quand la bibliothèque change. On s'en aperçoit au pire moment, la
veille d'un envoi. L'\emph{intégration continue} consiste à confier ces vérifications à une
machine, qui les refait à chaque modification envoyée sur le dépôt distant.

\section{Le principe}

On écrit dans le dépôt un fichier qui décrit les vérifications. À chaque envoi
(\cmd{git push}), le service d'hébergement (GitHub, GitLab) démarre une machine virtuelle, y copie le dépôt,
exécute les vérifications et affiche le résultat : une coche verte ou une croix rouge à côté
du commit. Sur GitHub, ce service s'appelle \emph{GitHub Actions} et le fichier vit dans
\fichier{.github/workflows/}.

\begin{terminal}
name: verifications
on: [push]
jobs:
  build:
    runs-on: ubuntu-latest
    steps:
      - uses: actions/checkout@v4
      - name: Compiler les preuves
        run: lake build
      - name: Compiler l'article
        run: tectonic -X compile article.tex
\end{terminal}

\section{Que vérifier ?}

Tout ce qui peut échouer sans qu'on le voie :
\begin{itemize}
  \item les preuves formelles, compilées en entier ;
  \item les documents \LaTeX{}, compilés, avec les avertissements de débordement de ligne
        traités comme des erreurs si l'on y tient ;
  \item le style du code (un \emph{linter}), et les types pour les programmes qui en ont ;
  \item les vérifications propres au projet : qu'une page citée existe, que chaque
        référence interne pointe quelque part, qu'un fichier de données respecte son schéma.
\end{itemize}

\begin{regle}
Un avertissement qu'on ne lit jamais est une erreur qu'on n'a pas encore vue. On fait
échouer la compilation sur les avertissements qui comptent.
\end{regle}

\section{Des tests pour les énoncés}

Les développeurs écrivent des tests avant ou pendant le code (le \emph{test-driven
development}) : un test dit ce que le programme doit faire, et l'intégration continue le
rejoue à chaque modification. Un projet de formalisation a ses tests tout trouvés. Chaque
énoncé Lean, avec \texttt{sorry} comme preuve, est un test qui échoue ; la preuve le fait
passer ; et l'intégration continue vérifie qu'aucun ne repasse au rouge quand la
bibliothèque change. On écrit les énoncés d'abord, comme on écrirait les tests.

\section{Le coût d'une erreur}

Une erreur vue à la minute coûte une minute. Vue par un rapporteur six mois plus tard, elle
coûte une révision et un peu de confiance. L'intégration continue ne trouve aucune erreur
mathématique. Elle attrape le reste, et l'attention reste libre pour les vraies.

\begin{tipblock}
Avant d'envoyer, on lance localement les mêmes commandes que la machine. Une vérification
locale manquée et rattrapée par l'intégration continue n'est pas grave ; mais si la croix
rouge devient habituelle, plus personne ne la regarde.
\end{tipblock}

\section{Déployer depuis la branche \texttt{main}}

Le même mécanisme publie. Quand les vérifications passent, une dernière étape construit le
site du projet et le met en ligne. La règle devient alors : ce qui est sur la branche
principale est en ligne, et ce qui est en ligne a passé les vérifications.
